% Table S16: P1 context used for P2
\begin{table}[htbp]
\centering
\scriptsize
\caption{Use of companion P1 evidence in P2. P1 results are included only as methodological or chemical-space context. They do not validate P2 mutant-state RRS, biological resistance, or polypharmacology.}
\label{tab:s16_p1_context}
\begin{tabularx}{\textwidth}{>{\raggedright\arraybackslash}p{0.22\textwidth} >{\raggedright\arraybackslash}p{0.30\textwidth} X}
\toprule
P1 evidence & Permitted use in P2 & Boundary \\
\midrule
Redocking and enrichment benchmarks & Context for pose reproduction and general ligand-ranking behavior of the docking workflow & Does not validate mutant-state score retention or RRS \\
Physicochemical and drug-likeness profiles & Context for describing the upstream chemical library and the selected PP-01--PP-17 cohort; the cohort audit found 17/17 exact SMILES matches & Shared candidate identity and provenance preclude independent-replication claims; properties are not biological activity measurements \\
Chemical-space and scaffold analyses & Context for the origin and diversity of the parent library & Set-C remains a filtered prioritization cohort, not a representative sample \\
MPO sensitivity & Context for the upstream selection procedure and its local stability & Does not test RRS or network-weighted ranking \\
P1--P2 shared candidates or provenance & Used only with explicit cohort and provenance labeling & Not an independent replication when candidates or source workflows overlap \\
\bottomrule
\end{tabularx}
\end{table}
